\section{Introduction}
\label{sec:intro}

Post-training is usually evaluated through behavior, but it also leaves behind a
geometric record: the difference between the initial and final weights. For a
linear map, this checkpoint delta is
$\dW=W_{\text{instruct}}-W_{\text{base}}$. It aggregates every stage between the
two checkpoints, so it cannot by itself distinguish instruction tuning,
preference optimization, or safety training
\citep{ouyang2022instruct,rafailov2023dpo}. It can, however, reveal where the
parameters moved. We ask whether that geometry can narrow the search for
directions associated with refusal and measured emergent misalignment.

The central difficulty is that geometric structure is not automatically
behavioral structure. A diffuse update spreads spectral mass across many
directions, whereas a concentrated update places substantial mass in a few.
Under an iid Gaussian diffuse-update model, the limiting spectrum has a
Marchenko--Pastur upper edge \citep{marchenko1967,johnstone2001}. Under an
additive fixed-rank model, a sufficiently strong component detaches from that
edge, and its singular vector partially overlaps the planted direction
\citep{benaych2012}. We use this theory only to decide which empirical
directions are spectrally visible. Behavioral meaning must come from separate
capture, control, and intervention experiments.

The paper follows that distinction from broad structure to specific behavior.
We first examine every linear map in the released Llama-3-8B Base-to-Instruct
delta \citep{grattafiori2024llama3}. All $224$ leading eigenvalues clear their
fitted Marchenko--Pastur visibility edges, often by one to four orders of
magnitude, and the median stable rank is $109$ against ambient dimensions in the
thousands. The update is therefore strongly anisotropic, but this census alone
does not make the anisotropy specific to alignment.

We next ask whether the visible directions intersect a known behavioral axis.
The empirical refusal direction, estimated independently from harmful and
harmless prompts \citep{arditi2024}, is enriched in the leading singular
subspace of the attention-output delta. Projecting the top-$128$ subspace out
after every block changes substring-scored refusal on a fixed harmful-prompt
slice from $98.4\%$ ($[94.5,99.6]\%$) to $14.1\%$ ($[9.1,21.1]\%$), while one
seeded random subspace gives $97.7\%$ ($[93.3,99.2]\%$)
\citep{wilson1927}. The same projection damages MMLU, ARC-C, and GSM8K far more
than the random control. The result links the leading subspace to measured
refusal, but the intervention is blunt and does not preserve capability.

Finally, we replace the unrestricted Base-to-Instruct delta with a matched
contrast. In the emergent-misalignment model-organism setting
\citep{betley2025emergent,turner2025organisms}, harmful and safe arms share
questions and training recipe while differing in target-answer policy and
content. Their difference yields a medical-contrast direction. Across four
Qwen2.5-Coder-7B-Instruct pairs, seed-level directions agree with the pooled
direction at mean cosine $0.97$, compared with $0.16$ for benign-pair
differences at the selected layer. A retrospective leave-one-pair-out score
ranks every harmful arm above its safe counterpart. In an in-sample
intervention, projecting out the direction changes the all-output joint
misalignment-and-eligibility rate from $2.3\%$ ($18/800$) to $0.0\%$ ($0/800$),
versus $3.4\%$ ($27/800$) under random-direction ablation. The pattern recurs
within three 7B/8B medical organisms, although the Mistral effect is partial;
the 14B causal audit does not meet its frozen positive criteria. A separate
$16$-pair Qwen audit shows that weight-SVD, row-mean contrast, and arm-labeled
activation PCA all separate held-out arms, but does not support superiority of
one weight method.

This sequence defines the scope of the contribution. The additive-spike theory
provides a visibility rule, the full-matrix census identifies candidate
subspaces, and matched contrasts and interventions test their behavioral
relevance. The results support behavior-sensitive geometry in the fixed
checkpoints studied here. They do not identify the training mechanism, establish
one-dimensional sufficiency, or demonstrate prospective causal generalization.
Figure~\ref{fig:pipeline} summarizes the first stage of the analysis.
